\documentclass[12pt]{letter}
\usepackage[margin=1.2in]{geometry}
\usepackage{hyperref}

\address{Independent Researcher\\Nanjing, China}
\signature{Huang Zhongchang}

\begin{document}

\begin{letter}{Editor\\Physical Review D}

\opening{Dear Editor,}

We submit ``A Quantum Many-Body Theory of Information Dark Energy'' for consideration in Physical Review D.

\textbf{What this paper does.} Gough [Entropy 27, 110 (2025)] recently identified an empirical correlation between the cosmic star formation history and the dark energy equation of state favored by DESI DR2---but no microscopic theory was available. We provide that theory. Dark energy, we propose, emerges from a driven-dissipative ensemble of cosmic two-level systems (qubits), each representing a quantum of information capacity. Gravitational collapse processes information by determining qubits from $|0\rangle$ to $|1\rangle$. The collective dynamics follow a complex Langevin equation with mean-field nonlinearity; the phantom crossing ($w$ crossing $-1$) occurs when the determined fraction exceeds a critical value $a_c^2$, derived from the Curie-Weiss mean-field theory of the $N$-qubit transverse Ising model.

\textbf{What is new.} (1) A microscopic theory for Gough's empirical finding---the first to provide dynamical equations, a Hamiltonian, and derived (not fitted) parameters. (2) The phantom crossing mechanism as an energy-gap sign flip in the qubit ensemble---conceptually distinct from prior mechanisms (Hu 2005, quintom, DM-DE interaction). (3) The driving function computed from gravitational collapse power (zero phenomenological inputs), removing the SFR-dependence of Gough's treatment. (4) A proposed analog quantum simulation on 3--5 qubit chips as a future test.

\textbf{Why PRD.} The paper provides a new theoretical framework (quantum many-body theory applied to dark energy) with honest scoping---it does not claim a detection (evidence is $\sim 2\sigma$, indicative), it explicitly discusses the weakest assumption (Landauer's principle for self-gravitating systems), and it engages thoroughly with prior art (Gough 2025, Hu 2005, Martineau \& Brandenberger 2005, Zhao et al.\ 2009). We believe this combination of theoretical novelty and scientific honesty is appropriate for PRD.

\textbf{Relation to prior work.} This paper is explicitly positioned as the microscopic theory for Gough's empirical finding, not as a paradigm invention. All relevant prior art is cited and distinguished.

Suggested referees: Eric Linder (dark energy phenomenology), Wayne Hu (phantom crossing theory), and an expert in quantum many-body theory (Dicke model/Ising model).

\closing{Sincerely,}

\end{letter}
\end{document}
